Στις σύγχρονες αγωνιστικές διοργανώσεις, τα ελαστικά αποτελούν ίσως τον πιο καθοριστικό 
παράγοντα επίδοσης και στρατηγικής. Τα ελαστικά λειτουργούν εντός ενός
στενού παραθύρου θερμοκρασίας (\textit{thermal window}): κάτω από ένα κατώτατο όριο η σύνθεση του
υλικού δεν μπορεί να αποδώσει τη μέγιστη πρόσφυση, ενώ πάνω από ένα ανώτατο η αυξημένη
θερμοκρασία επιταχύνει τη φθορά του πέλματος και μειώνει την πρόσφυση. Η κατανόηση και η ακριβής
μοντελοποίηση αυτής της συμπεριφοράς έχει τεράστια πρακτική αξία τόσο για τη
βελτιστοποίηση της απόδοσης ανά γύρο όσο και για τον σχεδιασμό της στρατηγικής του αγώνα, καθώς η φθορά του πέλματος αποτελεί τον πρωταρχικό παράγοντα που ορίζει τη στρατηγική
αντικατάστασης ελαστικών (\textit{pit-stop strategy}). Κάθε τύπος ελαστικού παρουσιάζει διαφορετική 
ισορροπία μεταξύ επίδοσης και
ανθεκτικότητας, και η βέλτιστη επιλογή εξαρτάται από χαρακτηριστικά πίστας, συνθήκες αγώνα
και τη συμπεριφορά οδήγησης \cite{West2022, Tremlett2016}. Η ακριβής πρόβλεψη του ρυθμού φθοράς
για δεδομένο προφίλ οδήγησης είναι εξαιρετικά δύσκολη και συνεχίζει να αποτελεί
αντικείμενο έρευνας.

Οι προσομοιωτές χρόνου γύρου (lap time simulators) αποτελούν βασικό εργαλείο στην ανάπτυξη
αγωνιστικών οχημάτων. Η πλειονότητά τους βασίζεται στην προσέγγιση ψευδο-μόνιμης κατάστασης
(\textit{Quasi-Static Simulation -- QSS}), η οποία χρησιμοποιεί ένα διάγραμμα επιδόσεων (GGV
diagram) ως στατική χαρτογράφηση των δυνατοτήτων του οχήματος σε κάθε ταχύτητα. Η
χαρτογράφηση αυτή κατασκευάζεται συνήθως για σταθερές, ιδανικές συνθήκες ελαστικού --
σταθερή θερμοκρασία και μηδενική φθορά -- αγνοώντας τη δυναμική εξέλιξη της κατάστασης των
ελαστικών κατά τη διάρκεια του αγώνα.

Με βάση τις ανάγκες της σύγχρονης μορφής του μηχανοκίνητου αθλητισμού, αυτή η παραδοχή 
παραλείπει μεγάλο μερίδιο των συνθηκών του οχήματος και λόγω της ευαισθησίας των ελαστικών
εισάγει σημαντικό σφάλμα στα σενάρια όπου η κατάσταση τους
αποκλίνει από την ιδανική: κατά τους πρώτους γύρους προθέρμανσης (out-lap), σε συνθήκες υψηλής φόρτισης
που οδηγούν σε υπερθέρμανση, ή κοντά στο τέλος της ζωής των ελαστικών. Επομένως, η πρόβλεψη χρόνου γύρου χωρίς να λαμβάνεται υπόψη η θερμοκρασία των ελαστικών
αφενός μπορεί να διαφέρει σημαντικά από την πραγματικότητα, αλλά σημαντικότερα δεν επιτρέπει την 
ανάλυση της στρατηγικής ενός αγώνα προκαταβολικά ή σε πραγματικό χρόνο.

Στη βιβλιογραφία υπάρχει καταγεγραμένη έρευνα τα τελευταία χρόνια για την ανάπτυξη μοντέλων πρόβλεψης 
θερμοκρασίας και φθοράς -- είτε απλούστερα είτε πιο λεπτομερή όπου στην πλειοψηφία τους
χρησιμοποιούν μη-μόνιμους προσομοιωτές χρόνου γύρου (\textit{Transient lap time simulators}). Οι μη-μόνιμοι προσομοιωτές είναι συνήθως ακριβέστεροι έναντι των προσομοιωτών ψευδο-μόνιμης κατάστασης, ωστόσο, μπορεί να εισάγουν αβεβαιότητα
λόγω της ανάγκης τους για μοντέλο οδηγού, και επίσης είναι πιο ασταθείς και υπολογιστικά κοστοβόροι. Οι προσομοιωτές ψευδο-μόνιμης κατάστασης χρησιμοποιούνται
κυρίως σε μελέτες αρχιτεκτονικής του οχήματος όπου απαιτείται η χαρτογράφηση των μέγιστων δυνατών επιδόσεων του οχήματος.

\subsection*{Σκοπός και συνεισφορά}

Ο σκοπός της παρούσας διπλωματικής εργασίας είναι η ανάπτυξη ενός υπολογιστικού πλαισίου προσομοίωσης χρόνου γύρου ψευδό-μόνιμης κατάστασης λαμβάνοντας υπόψη
την επίδραση της θερμοκρασίας και της φθοράς. Συγκεκριμένα, αναπτύσσεται ένα 
μοντέλο που μοντελοποιεί τη
θερμοκρασία και τη φθορά των ελαστικών κατά τη διάρκεια ενός γύρου, και ενσωματώνει αυτή
τη δυναμική στον υπολογισμό του διαγράμματος επιδόσεων - GGV. Επιλέγεται η ανάπτυξη ενός φυσικού θερμοδυναμικού μοντέλου έναντι άλλων 
(όπως για παράδειγμα μοντέλου μηχανικής μάθησης \cite{Todd_2025}) αφενός γιατί επιτρέπει την κατανόηση των μηχανισμών των φαινομένων που λαμβάνουν χώρα και
αφετέρου, μετά την επαλήθευση του μοντέλου, μπορεί εύκολα να γενικευθεί σε περισσότερους συνδυασμούς οχήματος--ελαστικών--πίστας.
Το αποτέλεσμα του μοντέλου είναι ένα
\emph{χρονο-μεταβαλλόμενο} διάγραμμα GGV που αντικατοπτρίζει τη στιγμιαία κατάσταση των ελαστικών,
και συνδέεται με τον προσομοιωτή QSS ανοιχτού κώδικα \textit{openLAP} για την εκτίμηση χρόνου γύρου.

Η κύρια συνεισφορά της εργασίας συνοψίζεται:
\begin{enumerate}
    \item Ανάπτυξη φυσικού θερμικού μοντέλου ελαστικού.
    \item Μοντελοποίηση τριών ανεξάρτητων μηχανισμών φθοράς.
    \item Σύνδεση θερμοκρασίας--πρόσφυσης για τη δυναμική ανανέωση του GGV.
    \item Διαδικασία υπολογισμού θερμοκρασίας: χρησιμοποιώντας δεδομένα πίστας/ταχύτητας απο προσομοίωση ή πειραματικά δεδομένα 
    (εφαρμογή στην πίστα της Βαρκελώνης -- Circuit de Barcelona-Catalunya) 
    \item Προσομοίωση χρόνου γύρου λαμβάνοντας υπόψη τη θερμοκρασία και φθορά των ελαστικών
    \item Μελέτη διατήρησης ελαστικών -- συσχέτιση χρόνου γύρου με συνολική φθορά% [TODO]
\end{enumerate}

\subsection*{Δομή εργασίας}

Στη συνέχεια, το Κεφάλαιο \cref{sec:tyre_theory} ξεκινά παρουσιάζοντας τη θεωρητική βάση για τη
συμπεριφορά ελαστικών: μηχανισμούς παραγωγής θερμότητας, απαγωγής και φθοράς. Το
Κεφάλαιο \cref{sec:tyre_model} αναπτύσσει τα υπολογιστικά μοντέλα πρόσφυσης, θερμότητας και
φθοράς με βάση τη θεωρία που παρουσιάστηκε. Έπειτα το Κεφάλαιο \cref{sec:vehicle_model} περιγράφει 
το μαθηματικό μοντέλο του οχήματος. [TODO] Να ολοκληρώσω την παράγραφο αφού τελειώσουν ολα τα κεφάλαια.
